% ---------------------------------------------------------------------------
% Appendix: Experimental details
% Requires: booktabs, amsmath (both standard in the ICLR template).
% Hyper-parameters are read from adversarial_demo/config.yaml, the merged run
% configs (results/cluster_eval*/<ds>_run_config_merged.yaml) and
% scripts/cluster/config.sh; timings from /Data/mathias.ollu/tmp/attack_timings.json.
% ---------------------------------------------------------------------------
\section{Experimental details}
\label{app:experimental-details}

\paragraph{Attacked models.}
Unless stated otherwise the attacked model is the public RFM~$\mathcal{S}^2$ checkpoint of
\citet{dufour2024around} for each dataset (\texttt{PLONK\_YFCC} and \texttt{PLONK\_OSV\_5M},
$\approx$36M parameters in the velocity network~$\phi$). The image encoder~$\mathcal{E}$ is
frozen and shared with the Diffusion and Euclidean Flow Matching checkpoints used in the
transfer experiment of Figure~\ref{fig:transfer}: DINOv2 ViT-L/14 with registers
\citep{oquab2024dinov2} on YFCC4k and the StreetCLIP ViT-L/14 vision tower
\citep{haas2023streetclip} (CLS token) on OSV-5M. Both encoders take a $336{\times}336$ input
obtained by a $1{:}1$ centre crop followed by bicubic resizing, and normalise it with the
ImageNet (DINOv2) or CLIP (StreetCLIP) channel statistics.

\paragraph{Data.}
YFCC4k is used in full ($4\,000$ images). For OSV-5M we evaluate a fixed subset of $5\,000$
test images obtained by a seeded shuffle (seed~42) of the test split; the selection is
prefix-stable so that every experiment sees the same images.
% NOTE(mo): Section 6.1 currently says 4 000 OSV images; the run used 5 000
% (osv_restarts_results.json: n_images = 5000). Align one or the other.

\paragraph{Shared attack protocol.}
Every attack of Table~\ref{tab:framework} is trained per image with the same outer loop
(Table~\ref{tab:hyperparams-shared}). The perturbation $\delta$ lives in the encoder input
space and is constrained to the $\ell_\infty$ ball of radius
$b\in\{4/255,\,8/255\}$ by clamping after every step. All denoiser-based attacks use
sign-SGD (PGD) with step size $10^{-3}$ for $60$ steps and are initialised at
$\delta_0\sim\mathcal{U}(-10^{-4},10^{-4})$; GeoFA-Encoder uses Adam with the same learning
rate and a random start $\delta_0\sim\mathcal{U}(-b,b)$. Expectations over $(\epsilon,t)$ are
estimated with a Monte-Carlo batch of $256$ samples per step, $t\sim\mathcal{U}[0,1]$. The data
endpoint $x_0$ needed by the forward process is not the ground-truth location: for every image
we draw once a bank of $512$ locations from the \emph{clean} model ($100$ sampling steps,
$w=0$) and resample $x_0$ from this bank at every step, so that no training-time attack
(GeoFA-DPA excepted, see below) requires the true coordinates. The bank is shared by all
restarts of an image.

\paragraph{Restarts and selection.}
Each attack (GeoShield excepted) is trained from $8$ independent random initialisations.
The reported perturbation is the best of the first $5$ restarts (Section~\ref{sec:setup});
restarts $6$--$8$ are only used for the restart ablation. A restart is scored by running the
sampler on the clean and on the protected image from the \emph{same} $128$ initial noise
vectors ($250$ steps, $w=10$, seed $42{+}\text{restart index}$) and measuring the mean
haversine distance between the two sets of final predictions; the restart with the largest
displacement is kept.
% NOTE(mo): selection uses displacement from the *clean prediction*
% (restart_selection_metric = final_step_displacement), not FSD w.r.t. the ground truth
% as Section 6.1 currently states. Either wording or code should change.

\begin{table}[h]
\centering\small
\caption{Hyper-parameters shared by all GeoFA attacks. Values are identical on YFCC4k and
OSV-5M.}
\label{tab:hyperparams-shared}
\begin{tabular}{ll}
\toprule
Budget $b$ ($\ell_\infty$)                         & $4/255$, $8/255$ \\
PGD steps                                          & $60$ \\
Step size / learning rate                          & $10^{-3}$ \\
Optimiser                                          & sign-SGD (Adam for GeoFA-Encoder) \\
Initialisation $\delta_0$                          & $\mathcal{U}(-10^{-4},10^{-4})$
                                                     ($\mathcal{U}(-b,b)$ for GeoFA-Encoder) \\
Monte-Carlo batch ($\epsilon$, $t$, $x_0$ per step) & $256$ \\
$x_0$ bank (clean-model samples)                   & $512$ samples, $100$ steps, $w=0$ \\
Restarts trained / reported                        & $8$ / best of $5$ \\
Restart scoring                                    & $128$ shared-noise samples, $250$ steps, $w=10$ \\
Seed                                               & $42$ \\
\bottomrule
\end{tabular}
\end{table}

\paragraph{Attack-specific settings.}
Table~\ref{tab:hyperparams-specific} lists the choices that differ between attacks.
\emph{GeoFA-Encoder} maximises the $\ell_2$ distance between the perturbed and clean
embeddings (untargeted). \emph{GeoFA-Sampling} backpropagates through the full $16$-step
RFM sampler run at $w=0$ (guidance would double the graph); the clean and perturbed samplers
share their initial noise and the loss is the mean haversine distance between the two sets of
$256$ endpoints. \emph{GeoFA-AdvDM} maximises the RFM training loss of
equation~\ref{equ:rfm-loss} with $x_t$ and $v(x_t)$ obtained from the geodesic forward process
between a uniform point on $\mathcal{S}^2$ and $x_0$. \emph{GeoFA-UniDef} optimises the same
objective but replaces autograd through $\phi$ by the finite-difference estimator of
equation~\ref{equ:unidef}: one probe direction $z=\mathcal{E}(c)/\lVert\mathcal{E}(c)\rVert$
and a symmetric step $h=10^{-2}$, i.e. two extra forward passes of $\phi$ and none backward.
\emph{GeoFA-ACE} uses the MIST target image of \citet{zheng2025ace} as $c'$ and weights the
encoder term with $\alpha=1$. \emph{GeoFA-DPA} steers towards the antipode of the image's
ground-truth coordinates, $x_0'=-x_0$, by minimising the training loss of the geodesic path
ending at $x_0'$. \emph{GeoFA-$L^2$} and \emph{GeoFA-TD} compare the perturbed velocity to the
clean velocity $\psi(x_t\mid c)$ recomputed at every step on the same $(x_t,t)$, with the
(negative) squared $\ell_2$ distance and the cosine similarity respectively.
\emph{GeoShield} is run with the authors' implementation and default settings:
$\epsilon\in\{4,8\}$ on the $0$--$255$ scale, $100$ FGSM steps of size $1/255$, an ensemble of
CLIP ViT-B/16, ViT-B/32 and LAION ViT-B/32 image encoders, $640$-pixel inputs and random
crops with scale in $[0.5,0.9]$; it has no restart mechanism and its perturbations are
produced directly in pixel space.

\begin{table}[h]
\centering\small\setlength{\tabcolsep}{5pt}
\caption{Attack-specific hyper-parameters and cost of one PGD step. ``fwd/bwd'' counts
forward and backward passes through the encoder $\mathcal{E}$ (one image) and the velocity
network $\phi$ (batch of $256$). GeoFA-Sampling unrolls $16$ sampler steps.}
\label{tab:hyperparams-specific}
\begin{tabular}{lllcc}
\toprule
Attack & Reference / target & Specific settings & $\mathcal{E}$ fwd/bwd & $\phi$ fwd/bwd \\
\midrule
GeoFA-Encoder  & clean embedding      & Adam, random start          & 1/1 & 0/0 \\
GeoFA-Sampling & clean prediction     & $16$ steps, $w=0$, shared noise & 1/1 & 32/16 \\
GeoFA-AdvDM    & training target      & --                          & 1/1 & 1/1 \\
GeoFA-UniDef   & training target      & FDJE, $h=10^{-2}$, $1$ probe & 1/1 & 3/0 \\
GeoFA-ACE      & MIST image           & $\alpha=1$, $\ell_2$ targets & 2/1 & 2/1 \\
GeoFA-DPA      & antipode of $x_0$    & requires true GPS           & 1/1 & 1/1 \\
GeoFA-$L^2$    & clean velocity       & $-\lVert\cdot\rVert^2$      & 2/1 & 2/1 \\
GeoFA-TD       & clean velocity       & cosine similarity           & 2/1 & 2/1 \\
\midrule
GeoShield      & clean geo-features   & $100$ steps, $\alpha=1/255$, CLIP ensemble & -- & -- \\
\bottomrule
\end{tabular}
\end{table}

\paragraph{Evaluation protocol.}
For every (image, attack, budget) the protected image is fed to the model exactly like a
clean one (crop, resize, normalisation). We draw $128$ locations with $250$ sampling steps
and guidance $w=10$ from initial noise shared with the clean run; FSD is the mean haversine
distance between these $128$ final predictions and the ground truth, and
$\mathrm{ASR}_{>d}$ is the fraction of images whose FSD exceeds $d\in\{200,750,2500\}$\,km.
The ``Clean'' rows of Table~\ref{tab:main} apply the same protocol to the unprotected image.
% NOTE(mo): see the message accompanying this file -- the current code passes the protected
% image through utils.adversarial_utils.tensor_to_pil, whose inverse normalisation does not
% match the encoder preprocessing; this sentence describes the intended protocol.

\paragraph{Ablations.}
All ablations reuse the selected perturbations; nothing is retrained.
\emph{Sampling steps:} the protected image is re-evaluated with
$\{8,16,32,64,250\}$ steps ($w=10$). \emph{Robustness:} the protected image is
JPEG-compressed at quality $\{20,40,60,80\}$ (100 denotes no compression) or blurred with a
Gaussian kernel of $\sigma\in\{0,2,4,6,8,10\}$ pixels before the encoder, then evaluated at
the default $250$ steps; this sweep is run for GeoShield, GeoFA-Encoder, GeoFA-Sampling and
GeoFA-TD. \emph{Transfer:} the same protected image is evaluated against the Diffusion
(DDIM) and Euclidean Flow Matching checkpoints of \citet{dufour2024around} that share the
encoder, at $250$ steps and $w=10$. \emph{Restarts:} success rates are recomputed when
keeping the best of the first $k\le 8$ restarts. \emph{Localizability:} the entropy of
equation~\ref{equ:localizability} is estimated on the clean image with $256$ importance
samples of the model likelihood, and images are split at the $33\%$ and $66\%$ quantiles.

\paragraph{Compute.}
Table~\ref{tab:runtime} reports wall-clock times per image measured on a single NVIDIA RTX
A5000 (24\,GB) with the YFCC4k model at $b=8/255$, averaged over three images. Training cost
is dominated by the number of passes through $\phi$ per step (Table~\ref{tab:hyperparams-specific});
the $128$-sample, $250$-step evaluation of a restart costs $1.7$\,s and the $x_0$ bank
$0.7$\,s. With the full protocol ($5$ restarts, each scored) protecting one image takes
between $45$\,s (GeoFA-Encoder) and $91$\,s (GeoFA-Sampling); GeoFA-TD takes $69$\,s, the same
as GeoFA-$L^2$ and $30\%$ more than GeoFA-AdvDM because it also evaluates the clean velocity.
The full-scale runs were sharded into jobs of $100$ images on a SLURM cluster; the seven
in-process attacks of Table~\ref{tab:main} amount to roughly $1.6$k (YFCC4k) and $2.0$k
(OSV-5M) A5000-GPU-hours including the three extra restarts, plus the (cheap) ablation
re-evaluations and the out-of-process GeoShield runs.
% NOTE(mo): OSV timings not measured (StreetCLIP encoder); GeoShield not timed in our harness.

\begin{table}[h]
\centering\small
\caption{Per-image run time (seconds, RTX A5000, YFCC4k model, $b=8/255$). ``Restart'' is
one $60$-step training run; ``Image'' is the full protocol: $x_0$ bank (when used)
$+\,5\times$(training + scoring of a restart).}
\label{tab:runtime}
\begin{tabular}{lccc}
\toprule
Attack & Training / restart & Scoring / restart & Total / image \\
\midrule
GeoFA-Encoder  &  7.4 & 1.7 & 45 \\
GeoFA-Sampling & 16.6 & 1.7 & 91 \\
GeoFA-AdvDM    &  8.6 & 1.7 & 52 \\
GeoFA-UniDef   &  8.8 & 1.7 & 53 \\
GeoFA-ACE      & 12.2 & 1.7 & 70 \\
GeoFA-DPA      &  8.6 & 1.7 & 52 \\
GeoFA-$L^2$    & 12.0 & 1.7 & 69 \\
GeoFA-TD       & 12.0 & 1.7 & 69 \\
\bottomrule
\end{tabular}
\end{table}
